\begin{longtable}{@{}p{0.6cm}p{7.2cm}p{2.0cm}p{3.2cm}@{}}
\caption{Planned future work for the major project (owners are proposed from the stated roles; status is as at the end of the minor project)}\label{tab:future}\\
\toprule \# & Planned work & Proposed lead & Status\\ \midrule \endfirsthead
\toprule \# & Planned work & Proposed lead & Status\\ \midrule \endhead
1 & Develop a complete QSFD simulation application that integrates fraud detection, cryptographic security, quantum security experiments, attack simulation and result visualization in one platform. & Whole team & Not started (components exist as scripts)\\
2 & Create a real-time banking transaction simulator in which normal and fraudulent transactions can be generated or entered manually and passed through the complete QSFD pipeline. & Shreyansh & Not started\\
3 & Implement comparative fraud detection using Logistic Regression, Random Forest, XGBoost, QSVM, VQC, QNN and the final quantum-enhanced ensemble. & Shreyansh & Partly done: models compared in Layer 1; unified comparison inside the application remains\\
4 & Develop a dynamic transaction risk-scoring system that classifies transactions into Low, Medium and High risk from the predicted fraud probability and other transaction characteristics. & Shreyansh & Partly done: risk levels exist in the policy engine\\
5 & Implement risk-adaptive security selection so that mechanisms are chosen according to transaction risk instead of applying the most expensive configuration to every transaction. & Shreyansh & Partly done: four security profiles and policies implemented\\
6 & Compare fixed security with risk-adaptive security to determine whether adaptive security reduces average latency and computational overhead while maintaining the intended protection. & Abhishek & Partly done: routing ablation completed; latency comparison pending\\
7 & Perform classical versus post-quantum cryptography experiments comparing RSA and ECC with ML-KEM, ML-DSA and SLH-DSA in execution time, key size, ciphertext size, signature size and verification time. & Abhishek & Done in the minor project; to be extended\\
8 & Compare reference and optimized PQC implementations to determine how much performance depends on the implementation backend rather than only on the algorithm. & Abhishek & Done in the minor project (pqcrypto vs pure Python); native liboqs still pending\\
9 & Perform complete end-to-end banking security experiments covering transaction creation, fraud detection, security selection, key establishment, signing, AES-256 encryption, transmission simulation, decryption, signature verification and the final transaction decision. & Abhishek and Debshata & Partly done: cryptographic pipeline measured; fraud-detection and decision stages not yet chained\\
10 & Perform scalability and load experiments with increasing numbers of simulated transactions (1, 10, 100, 1,000 and more), measuring latency, throughput, memory usage and security-processing overhead. & Shreyansh & Not started\\
11 & Perform transaction-tampering experiments by modifying amounts, recipient details, ciphertexts or digital signatures after protection and checking whether the security layer reliably detects the manipulation. & Abhishek & Not started (tamper rejection is unit-tested)\\
12 & Implement a replay-attack simulation in which a previously valid transaction is submitted again, to study whether transaction IDs, timestamps, nonces or similar mechanisms detect repetition. & Abhishek & Not started\\
13 & Expand the BB84 simulation laboratory with different channel-noise levels, key lengths, numbers of trials and Eve interception probabilities, measuring QBER and sifting efficiency. & Debshata & Partly done: noise levels and trials covered\\
14 & Perform partial-eavesdropping experiments in which Eve intercepts 0\%, 25\%, 50\%, 75\% and 100\% of the transmitted BB84 qubits, to determine how attack intensity affects QBER and detection probability. & Debshata & Not started (full interception done)\\
15 & Evaluate Eve detection scientifically by varying the QBER threshold and calculating true-positive, false-positive, false-negative and miss rates, instead of assuming that one threshold proves the presence of Eve. & Debshata & Done for the full-interception case in the minor project\\
16 & Extend IBM Quantum hardware validation by comparing theory, ideal Aer, backend-derived noisy Aer and repeated hardware execution, while studying physical qubit quality, readout error, T1/T2, transpilation, circuit depth and run-to-run variability. & Debshata & Done on one backend (ibm\_marrakesh); more backends and calibration windows pending\\
17 & Extend the Grover experiment with small search spaces and compare ideal simulation, noisy simulation and IBM hardware, to study how hardware noise affects search success, without claiming a practical attack on AES. & Debshata & Partly done: two-qubit Grover\\
18 & Implement a small educational Shor experiment if time permits, such as factoring 15, strictly as a demonstration of the principle behind the future quantum threat to RSA and not as evidence that practical RSA keys can currently be broken. & Debshata & Not started\\
19 & Develop an interactive research dashboard (for example with Streamlit) showing transaction details, fraud probability, model predictions, selected security mechanism, cryptographic latency, payload overhead, BB84 QBER, Eve detection, IBM hardware results, Grover results and comparison graphs. & Shreyansh (with the team) & Not started\\
20 & Conduct a final integrated QSFD experiment on a synthetic stream of banking transactions, measuring the complete system from arrival to fraud decision, security selection, protection, attack/tampering detection, approval or blocking, and result storage; this becomes the major project's principal system-level experiment. & Whole team & Not started\\
\bottomrule
\end{longtable}
